% =============================================================================
% Anexo B · Robustez con las series transformadas (decisión D15)
% Fuentes: data/results/final/anexo/ (05c_*.csv y 09b_discriminacion_episodios.csv),
% notebooks 05c y 09b, R/ciclo_financiero_final.Rmd (anexo R). Cifras solo por macro.
% =============================================================================
\chapter{Robustez con las series transformadas}
\label{ap:robustez}

Por la decisión D15, el análisis principal usa las series sin transformar y las transformadas
pasan a este anexo con dos fines: sustentar la decisión de no diferenciar
(\cref{sec:ap-no-diferenciar}) y evaluar la detección de episodios fuera de muestra, donde solo la
vía transformada es informativa (\cref{sec:ap-deteccion}). Las series transformadas cubren
\nMesesTransformadas{} meses, de \fechaInicioTransformadas{} a \fechaFinDatos{}, con las mismas
\nSeries{} series del análisis principal. Este anexo no reabre la Parte~C de la decisión D2, que
sigue pendiente de discutir con el asesor (\cref{sec:pend-parte-c}).

\section{Sustento de no diferenciar}
\label{sec:ap-no-diferenciar}

La transformación de las series combina dos pasos: el logaritmo, en las series de magnitudes
positivas, y la primera diferencia. El notebook \texttt{05c} los separa y repite el PCA después
de cada uno.

\paragraph{El logaritmo no cambia la estructura.} Tras aplicar solo el logaritmo, la congruencia
de Tucker (\cref{def:tucker}) de los tres primeros ejes con los de las series sin transformar es
\congruenciaLogPCUno{},
\congruenciaLogPCDos{}
y \congruenciaLogPCTres{},
y la inercia del primer eje casi no cambia. La discusión sobre la transformación es, por tanto,
una discusión sobre diferenciar.

\paragraph{La diferencia sí la cambia.} Con solo la primera diferencia, la inercia del primer
eje cae de \inerciaPCUnoNiveles\,\% a
\inerciaPCUnoDiferencia{}\,\%,
y su correlación con el tiempo desaparece. La \cref{prop:ganancia} explica por qué: la primera
diferencia multiplica la varianza de una oscilación de período $P$ meses por
$\gamma(P)^2=4\sin^2(\pi/P)$, que es menor que uno para todo $P>6$ y tiende a cero como
$(2\pi/P)^2$. Una oscilación de ocho años ($P=96$) conserva solo la fracción $4\sin^2(\pi/96)$
de su varianza; el archivo \texttt{05c\_ganancia.csv} tabula la ganancia para los períodos de
interés. En los datos, la proporción de la varianza de cada serie en oscilaciones de dos años o
más pasa, en mediana, de
\espectroLentoMedianaNivelesSinTendencia\,\%
en las series sin transformar y sin tendencia lineal a
\espectroLentoMedianaTransformadas\,\%
en las transformadas.

\paragraph{Lectura.} Diferenciar elimina justamente la escala de tiempo que define el ciclo
financiero (\cref{sec:ea-ciclo}). Analizar las series sin transformar conserva esa escala, a
cambio de mezclar la tendencia de profundización con la estructura de régimen, mezcla que la
regla D5 y el nulo de tendencias mixto permiten leer (\cref{sec:res-tendencia}).

\section{Detección de episodios fuera de muestra}
\label{sec:ap-deteccion}

La razón de error \emph{walk-forward} del PCA (\cref{def:walkforward}) se calcula con $d=3$ en
las dos vías. Con las series transformadas, la proporción de falsas alarmas fuera de los
episodios ampliados en 3 meses es
\falsasAlarmasTransformadasTres{}\,\%,
y la sensibilidad en la pandemia de COVID-19 es
\sensibilidadCovidTransformadasTres{}\,\%.
Con las series sin transformar, en cambio, las falsas alarmas suben a
\falsasAlarmasNivelesTres{}\,\%:
el PCA ajustado con los meses anteriores no reconstruye bien los meses nuevos porque el nivel de
las series sigue creciendo, de modo que casi todo mes parece atípico. Por eso la decisión D15
conserva la detección de episodios en la vía transformada.

Dos limitaciones acotan esta lectura. Los episodios son pocos y algunos puntuales, de modo que la
sensibilidad se calcula sobre muy pocos meses y no tiene poder estadístico. Y los episodios que
la vía transformada detecta son choques abruptos, como la pandemia, no cambios lentos de
estructura, que son los que interesan a las hipótesis HT1 y HT3. La complementariedad entre ambas
vías es precisamente la propuesta de la Parte~C de D2, que se discutirá con el asesor.
